\subsection{How the strands inform each other}
\label{sec:integration}
The strands meet at interpretation. The survey describes how common particular reports were in the recruited sample and how they vary with other reported conditions. The migrant synthesis helps distinguish steps or decisions that the same survey response might conceal. The institutional interviews and regional experts describe implementation; official records identify formal responsibilities. Table~\ref{tab:joint-evidence} asks what each can add to the same substantive question without requiring them to measure the same outcome \citep{greeneCaracelliGraham1989,fettersCurryCreswell2013,guettermanFettersCreswell2015}.

We compare the unit, period, direction and stage of the transaction, and how the information was elicited. When sources speak to different stages, we treat them as complementary rather than as replications. Agreement requires a comparable claim; so does disagreement. When an item was not asked or a step is undocumented, we treat it as unobserved, not as a negative response \citep{ocathainMurphyNicholl2010}. An account repeated in a memo and a report is one source. We draw on \citet{humphreysJacobs2023} in asking which observations can distinguish plausible interpretations and which cannot. We do not link participant records, combine evidence in a likelihood, or infer a policy effect.

\begin{singlespace}
\small
\setlength{\tabcolsep}{3.5pt}
\begin{longtable}{@{}
  >{\raggedright\arraybackslash}p{0.138\linewidth}
  >{\raggedright\arraybackslash}p{0.203\linewidth}
  >{\raggedright\arraybackslash}p{0.199\linewidth}
  >{\raggedright\arraybackslash}p{0.180\linewidth}
  >{\raggedright\arraybackslash}p{0.213\linewidth}@{}}
\caption{Joint display: what the evidence supports about access, use and control}
\label{tab:joint-evidence}\\
\toprule
\textbf{Question} & \textbf{Survey evidence} & \textbf{Migrant synthesis} & \textbf{Institutions and experts} & \textbf{Joint reading and limit} \\
\midrule
\endfirsthead
\toprule
\textbf{Question} & \textbf{Survey evidence} & \textbf{Migrant synthesis} & \textbf{Institutions and experts} & \textbf{Joint reading and limit} \\
\midrule
\endhead
\midrule
\multicolumn{5}{r@{}}{\footnotesize Continued on next page}\\
\endfoot
\bottomrule
\endlastfoot
Domestic use and remittance journeys
& Domestic digital payments or transfers in 60 days: 368/413. Digital retention or use of incoming proceeds: 169/423. Bases are valid answers.
& Local payments can start a transfer to Venezuela. Some incoming transfers follow more direct routes.
& Domestic interoperation does not ensure conversion and delivery abroad.
& Different stages and directions. Use can coexist with dependence; these are not two rates for the same transfer. \\
\addlinespace
Entry and account control
& Both own-name conditions versus one: positive association with digital retention or use of incoming proceeds; limited exact support.
& A document may not be accepted. Account holders can help others pay without giving them control.
& Verification and continuity after a change of document are separate problems.
& Measurement is narrower than access. Neither eligibility nor control over spending is verified. \\
\addlinespace
Ability and practice
& Higher score for reported transfer ability and confidence in avoiding fraud: positive association with domestic use; limited exact support.
& Women describe learning through use and assisting others.
& Regional experts report varied prior experience; local interviews flag app barriers.
& Order unresolved. Prior skill, learning through use and help can coexist. No skills test or estimated training effect. \\
\addlinespace
Price, timing and delivery
& Fee visibility: last operation in 90 days, asked of affirmative 60-day users. Inbound timing and rate clarity are distinct.
& Women check rates, coordinate with intermediaries and confirm that recipients can use the funds.
& Experts distinguish domestic speed from conversion and final delivery.
& A fast step or visible fee does not measure total time or cost. No welfare estimate. \\
\addlinespace
Precautions and remedy
& A reported payment or remittance problem in 12 months opens the complaint item. First response and resolution require a complaint.
& Women describe checking suspicious messages and false confirmations. Trust does not ensure recovery.
& The domestic Bre-B complaint route need not cover a foreign leg.
& Reported problems are not verified fraud. A response is not resolution, and resolution need not mean recovery. \\
\addlinespace
Rules and actual use
& No assignment to a reform or measure of the rail used for a transaction.
& Wallet names do not identify a rail. A hypothetical Bre-B exercise cannot establish use.
& The benchmark compares formal designs; experts describe the work of implementation.
& Roles, not impacts. No inference about compliance, migrant uptake or relative country performance. \\
\addlinespace
Observed profiles
& Four combinations of two reported own-name conditions (phone line and address proof): 404 scorable pairs; 19 unscored.
& The earlier positions overlap and can change with the task; security cuts across them.
& Experts raise provider checks and product limits not measured by these conditions.
& The profiles show conditions, not types or control. No case linkage or adopted class assignments. \\
\end{longtable}
\end{singlespace}

The joint reading changes the interpretation of use in two ways. First, a digital transaction can depend on someone else's account or on an intermediary who completes a later step. The migrant synthesis makes these arrangements visible; the survey does not tell us how often they occur among domestic users. Second, a documented remedy can cover a different step from the one that failed. The institutional evidence helps locate the boundary of that remedy; it does not establish how a user's complaint ended. Seen together, the strands identify what a count of payments misses.

They also leave real alternatives open. Intermediation can provide a useful service as well as introduce a dependency. Reported ability could precede use or grow with it. And concern about active wallet use in a regional expert memo is not directly comparable to reported use of any domestic digital payment or transfer in this sample. For these questions we retain the alternatives rather than choose a direction of effect or force agreement. The technical supplement documents the provenance and limits of the comparisons.
